
\chapter{Quasiprimitive Permutation Groups and Small Rank}

\section{The Exercise}

\begin{definition}[\texorpdfstring{$n$}{n}-transitive permutation group]
    Let \(G \leq \operatorname{Sym}(\Omega)\), and let \(n\) be a positive integer with
    \(n \leq |\Omega|\). We say that \(G\) is \textbf{\(n\)-transitive} on \(\Omega\) if
    \(G\) acts transitively on the set of ordered \(n\)-tuples of pairwise distinct points:
    \[
        \Omega^{(n)}
        =
        \{(\omega_1,\ldots,\omega_n)\in \Omega^n
          \mid \omega_i\neq \omega_j \text{ whenever } i\neq j\}.
    \]
    Equivalently, for any two ordered \(n\)-tuples of pairwise distinct points
    \[
        (\omega_1,\ldots,\omega_n),
        \qquad
        (\omega_1',\ldots,\omega_n')
        \in \Omega^{(n)},
    \]
    there exists \(g\in G\) such that
    \[
        \omega_i^g=\omega_i'
        \qquad
        \text{for all } 1\leq i\leq n.
    \]
    \end{definition}

\begin{exercise}
Let \(G \leq \operatorname{Sym}(\Omega)\) be a quasiprimitive permutation group.
Characterize the cases where the point stabilizer \(G_\omega\) has exactly
\[
r=2,3,4,5
\]
orbits on \(\Omega\). Equivalently, characterize quasiprimitive permutation groups
of rank \(r\), for \(r=2,3,4,5\).

In particular, determine the best possible lower bound for \(r\) when \(G\) is of
holomorph simple type.
\end{exercise}

\begin{proof}
    Assume throughout that \(G\) is finite and transitive on \(\Omega\), and write
    \[
    r=\operatorname{rank}(G)
    =
    \#\{\text{\(G_\omega\)-orbits on }\Omega\}.
    \]
    We use the quasiprimitive O'Nan--Scott types
    \[
    \mathrm{HA},\ \mathrm{AS},\ \mathrm{HS},\ \mathrm{HC},\ \mathrm{SD},\ \mathrm{CD},\ \mathrm{TW},\ \mathrm{PA}.
    \]
    
    The answer, at the level of O'Nan--Scott type, is as follows:
    \medskip
    \begin{center}
    \begin{tabular}{c|p{0.8\textwidth}}
    \toprule
    $r$ & Possibilities \\ \midrule
    $2$ & $\mathrm{HA}$ or $\mathrm{AS}$; equivalently, $G$ is $2$-transitive. \\[2mm]
    $3$ & $\mathrm{HA}$ or $\mathrm{AS}$ of rank $3$, or $\mathrm{PA}$ on $\Delta^2$ with $2$-transitive component. \\[2mm]
    $4$ & $\mathrm{HA}$ or $\mathrm{AS}$ of rank $4$, or $\mathrm{HS}$ of rank $4$, or $\mathrm{PA}$ on $\Delta^3$ with $2$-transitive component. \\[2mm]
    $5$ & $\mathrm{HA}$ or $\mathrm{AS}$ of rank $5$, or $\mathrm{HS}$ of rank $5$, or $\mathrm{PA}$ on $\Delta^4$ with $2$-transitive component and set-transitive top group. \\
    \bottomrule
    \end{tabular}
    \end{center}
    \medskip
    More explicitly, in the product action case
    \[
    \Omega=\Delta^k,
    \]
    the rank is \(k+1\) precisely when the component group is \(2\)-transitive on
    \(\Delta\) and the top group is transitive on \(j\)-subsets of
    \(\{1,\ldots,k\}\) for each \(0\leq j\leq k\). In that case the \(G_\omega\)-orbits
    are exactly the weight classes
    \[
    \Omega_j=\{(\delta_1,\ldots,\delta_k)\in \Delta^k
    \mid |\{i:\delta_i\neq \delta\}|=j\},
    \qquad 0\leq j\leq k,
    \]
    where \(\omega=(\delta,\ldots,\delta)\).
    
    We now justify the table.
    
    First suppose \(r=2\). Then \(G_\omega\) has exactly two orbits on \(\Omega\),
    namely
    \[
    \{\omega\}
    \quad\text{and}\quad
    \Omega\setminus\{\omega\}.
    \]
    Thus \(G_\omega\) is transitive on \(\Omega\setminus\{\omega\}\), so \(G\) is
    \(2\)-transitive. Conversely, every \(2\)-transitive group has rank \(2\).
    By Burnside's theorem, a finite \(2\)-transitive permutation group is either of
    affine type or almost simple type. Hence the rank \(2\) quasiprimitive groups are
    precisely the \(2\)-transitive groups of type \(\mathrm{HA}\) or \(\mathrm{AS}\).
    
    Now let \(N\) be a minimal normal subgroup of \(G\). Since \(G\) is
    quasiprimitive, every nontrivial normal subgroup is transitive; in particular, \(N\)
    is transitive. By the quasiprimitive O'Nan--Scott theorem, \(G\) is one of the
    types listed above. We show which types can occur when \(r\leq 5\).
    
    Consider first a regular nonabelian normal subgroup
    \[
    N\cong T^m,
    \]
    where \(T\) is nonabelian simple. Identifying \(\Omega\) with \(N\), and taking
    \(\omega=1\), we have
    \[
    G_\omega\leq \operatorname{Aut}(N).
    \]
    Hence \(G_\omega\)-orbits preserve element orders.
    
    Since \(T\) is nonabelian simple, \(|T|\) has at least three distinct prime divisors.
    Indeed, if \(|T|\) involved at most two primes, then \(T\) would be solvable by
    Burnside's \(p^aq^b\)-theorem, impossible. Choose elements
    \[
    x,y,z\in T
    \]
    of distinct prime orders \(p,q,s\), respectively.
    
    If \(m=1\), then the elements
    \[
    1,\quad x,\quad y,\quad z
    \]
    lie in four distinct \(G_\omega\)-orbits. Therefore
    \[
    r\geq 4.
    \]
    This is the holomorph simple case. Thus a group of holomorph simple type cannot
    have rank \(2\) or \(3\).
    
    If \(m\geq 2\), then in \(T^m\) the elements
    \[
    1,\quad
    (x,1,\ldots,1),\quad
    (y,1,\ldots,1),\quad
    (z,1,\ldots,1),\quad
    (x,y,1,\ldots,1),\quad
    (x,z,1,\ldots,1)
    \]
    have pairwise distinct orders
    \[
    1,\quad p,\quad q,\quad s,\quad pq,\quad ps.
    \]
    Hence \(r\geq 6\). Thus regular compound nonabelian types, such as
    \(\mathrm{HC}\) and \(\mathrm{TW}\), do not occur for \(r=2,3,4,5\).
    
    Next consider diagonal type. In simple diagonal type with
    \[
    N=T^k,\qquad k\geq 3,
    \]
    points may be represented by tuples
    \[
    (a_1,\ldots,a_k)\in T^k
    \]
    modulo diagonal multiplication. The point stabilizer preserves the multiset of
    orders
    \[
    \left\{
    |a_i a_j^{-1}| \mid 1\leq i<j\leq k
    \right\}.
    \]
    Choose elements of two distinct odd prime orders \(q\) and \(s\) in \(T\). Then the
    following six points give six distinct \(G_\omega\)-orbits:
    \[
    (1,\ldots,1),
    \]
    \[
    (1,x,1,\ldots,1),\qquad
    (1,y,1,\ldots,1),\qquad
    (1,z,1,\ldots,1),
    \]
    where \(x,y,z\) have distinct prime orders, and
    \[
    (1,y,y^{-1},1,\ldots,1),
    \qquad
    (1,z,z^{-1},1,\ldots,1).
    \]
    Indeed, the corresponding multisets of pairwise quotient orders are distinct.
    Thus diagonal type with \(k\geq 3\) has rank at least \(6\). The same argument
    excludes compound diagonal type for \(r\leq 5\). The remaining diagonal case
    \(k=2\) is the holomorph simple situation already considered.
    
    It remains to discuss product action type. Suppose
    \[
    \Omega=\Delta^k,
    \qquad
    \omega=(\delta,\ldots,\delta),
    \]
    and \(G\) is contained in a product action wreath product
    \[
    L\wr P
    =
    L^k\rtimes P,
    \]
    where \(L\) acts transitively on \(\Delta\) and \(P\leq S_k\) is transitive on the
    set of coordinates.
    
    Define the weight of a point \(\alpha=(\alpha_1,\ldots,\alpha_k)\) by
    \[
    \operatorname{wt}(\alpha)
    =
    |\{i:\alpha_i\neq \delta\}|.
    \]
    The point stabilizer \(G_\omega\) preserves weight. Hence
    \[
    r\geq k+1.
    \]
    If \(L\) is \(2\)-transitive on \(\Delta\), then \(L_\delta\) is transitive on
    \(\Delta\setminus\{\delta\}\). If, in addition, the top group \(P\) is transitive on
    \(j\)-subsets of \(\{1,\ldots,k\}\) for every \(j\), then \(G_\omega\) is transitive
    on every weight class. Therefore
    \[
    r=k+1.
    \]
    
    Conversely, suppose \(r\leq 5\). If \(L\) is not \(2\)-transitive, then
    \(L_\delta\) has at least two nontrivial orbits, say \(A\) and \(B\), on
    \(\Delta\setminus\{\delta\}\). Then, already for \(k\geq 2\), the following six
    types give distinct \(G_\omega\)-invariant subsets:
    \[
    0,\qquad
    1A,\qquad
    1B,\qquad
    2AA,\qquad
    2AB,\qquad
    2BB.
    \]
    Thus \(r\geq 6\), contradiction. Hence the component action must be
    \(2\)-transitive.
    
    Similarly, if the top group is not transitive on \(j\)-subsets for some \(j\), then
    the weight \(j\) class splits into at least two \(G_\omega\)-orbits. For \(r\leq 5\)
    this forces the top group to be set-transitive on the relevant coordinate subsets.
    Consequently, in product action type and rank at most \(5\), we must have
    \[
    r=k+1.
    \]
    Thus the product action cases are exactly:
    \[
    k=2 \quad (r=3),\qquad
    k=3 \quad (r=4),\qquad
    k=4 \quad (r=5),
    \]
    with \(2\)-transitive component action, and with set-transitive top group.
    
    Combining these exclusions with the unrestricted affine and almost simple cases gives
    the table above.
    
    Finally, consider the special question about holomorph simple type. In this case
    \[
    T\leq G\leq T\rtimes \operatorname{Aut}(T),
    \]
    where \(T\) is nonabelian simple and acts regularly on \(\Omega\cong T\). Taking
    \(\omega=1\), we have
    \[
    G_\omega\leq \operatorname{Aut}(T).
    \]
    Since automorphisms preserve element orders and \(T\) has elements of at least three
    distinct prime orders, the four elements
    \[
    1,\quad x,\quad y,\quad z
    \]
    lie in four distinct \(G_\omega\)-orbits. Hence
    \[
    r\geq 4.
    \]
    Therefore, if \(G\) is of holomorph simple type, then
    \[
    \boxed{r>3.}
    \]
    
    This bound is sharp. Take \(T=A_5\) and
    \[
    G=A_5\rtimes \operatorname{Aut}(A_5).
    \]
    Since
    \[
    \operatorname{Aut}(A_5)\cong S_5,
    \]
    the point stabilizer acts on \(A_5\) by conjugation. The nonidentity elements of
    \(A_5\) have orders \(2,3,5\), and \(S_5\) is transitive on the elements of each of
    these orders. Thus the \(G_\omega\)-orbits on \(\Omega\cong A_5\) are exactly
    \[
    \{1\},
    \qquad
    \{x\in A_5: |x|=2\},
    \qquad
    \{x\in A_5: |x|=3\},
    \qquad
    \{x\in A_5: |x|=5\}.
    \]
    Hence
    \[
    \operatorname{rank}(G)=4.
    \]
    So the lower bound \(r\geq 4\) is best possible.
    \end{proof}

\section{Rank of a Transitive Permutation Group}

\begin{definition}[Rank]
Let \(G \leq \operatorname{Sym}(\Omega)\) be transitive. Then \(G\) acts on
\[
\Omega \times \Omega
=
\{(\omega_1,\omega_2) \mid \omega_1,\omega_2 \in \Omega\}
\]
by
\[
(\omega_1,\omega_2)^g=(\omega_1^g,\omega_2^g).
\]
The number of \(G\)-orbits on \(\Omega \times \Omega\) is called the
\textbf{rank} of \(G\) on \(\Omega\), and is denoted by
\[
\operatorname{rank}(G).
\]
\end{definition}

\begin{proposition}
Let \(G \leq \operatorname{Sym}(\Omega)\) be transitive, and fix
\(\omega \in \Omega\). Then
\[
\operatorname{rank}(G)
=
\#\{G_\omega\text{-orbits on }\Omega\}.
\]
In particular,
\[
\operatorname{rank}(G)\geq 2
\]
whenever \(|\Omega|>1\), and
\[
\operatorname{rank}(G)=2
\]
if and only if \(G\) is \(2\)-transitive on \(\Omega\).
\end{proposition}

\begin{proof}
Every \(G\)-orbit on \(\Omega \times \Omega\) meets the subset
\[
\{\omega\}\times \Omega.
\]
Indeed, for any pair \((\alpha,\beta)\), choose \(g \in G\) with
\(\alpha^g=\omega\). Then
\[
(\alpha,\beta)^g=(\omega,\beta^g).
\]
Two pairs \((\omega,\beta_1)\) and \((\omega,\beta_2)\) lie in the same
\(G\)-orbit if and only if \(\beta_1\) and \(\beta_2\) lie in the same
\(G_\omega\)-orbit. Hence the rank of \(G\) is the number of
\(G_\omega\)-orbits on \(\Omega\).

The singleton \(\{\omega\}\) is always a \(G_\omega\)-orbit. Thus, if
\(|\Omega|>1\), there is at least one further orbit, so
\(\operatorname{rank}(G)\geq 2\).

Finally, \(\operatorname{rank}(G)=2\) precisely when \(G_\omega\) is transitive on
\(\Omega \setminus \{\omega\}\), which is equivalent to \(G\) being
\(2\)-transitive.
\end{proof}

\begin{remark}
The \(G_\omega\)-orbits on \(\Omega\) are often called the \textbf{suborbits} of
\(G\). Thus the rank of a transitive permutation group is the number of suborbits.
\end{remark}

\section{Two-Transitive Groups}

\begin{theorem}[Burnside]
Let \(G \leq \operatorname{Sym}(\Omega)\) be a finite \(2\)-transitive permutation
group. Then \(G\) is either of affine type or almost simple type.
\end{theorem}

\begin{remark}
In O'Nan--Scott language, the conclusion says that a finite \(2\)-transitive group
is either of holomorph affine type, usually denoted HA, or of almost simple type,
usually denoted AS.
\end{remark}

\begin{proposition}
Let \(G \leq \operatorname{Sym}(\Omega)\) be \(2\)-transitive, and fix
\(\omega \in \Omega\). Then:
\begin{enumerate}[label=\textup{(\roman*)}]
    \item \(G_\omega\) is faithful and transitive on
    \(\Omega \setminus \{\omega\}\);
    \item \(|\Omega|-1\) divides \(|G_\omega|\);
    \item if \(N \lhd G\), then \(N_\omega=N\cap G_\omega\) is half-transitive on
    \(\Omega \setminus \{\omega\}\), i.e. all \(N_\omega\)-orbits on
    \(\Omega \setminus \{\omega\}\) have the same size.
\end{enumerate}
\end{proposition}

\begin{proof}
Since \(G\) is \(2\)-transitive, \(G_\omega\) is transitive on
\(\Omega \setminus \{\omega\}\). If an element of \(G_\omega\) fixes every point
of \(\Omega \setminus \{\omega\}\), then it fixes every point of \(\Omega\), and
hence is the identity. Thus the induced action is faithful.

By orbit-stabilizer,
\[
|\Omega|-1
=
|\omega^{G_\omega}|
\mid |G_\omega|.
\]

Finally, if \(N \lhd G\), then
\[
N_\omega=N\cap G_\omega \lhd G_\omega.
\]
Since \(G_\omega\) is transitive on \(\Omega \setminus \{\omega\}\), the orbits of
the normal subgroup \(N_\omega\) on this set all have the same size.
\end{proof}

\section{Normal Subgroups and Quasiprimitivity}

\begin{definition}[Half-transitive subgroup]
Let a group \(H\) act on a finite set \(\Lambda\). We say that \(H\) is
\textbf{half-transitive} on \(\Lambda\) if all \(H\)-orbits on \(\Lambda\) have the
same size.
\end{definition}

\begin{lemma}
Let \(G \leq \operatorname{Sym}(\Omega)\) be transitive, and let \(N \lhd G\).
Then the \(N\)-orbits on \(\Omega\) form a \(G\)-invariant block system. In
particular, all \(N\)-orbits have the same size.
\end{lemma}

\begin{proof}
Let \(\Delta_1,\Delta_2\) be two \(N\)-orbits. Choose
\(\delta_i \in \Delta_i\), for \(i=1,2\). Since \(G\) is transitive, there exists
\(g \in G\) such that
\[
\delta_1^g=\delta_2.
\]
Write
\[
\Delta_1=\{\delta_1^{x_1},\delta_1^{x_2},\ldots,\delta_1^{x_m}\},
\qquad x_i \in N.
\]
Then
\[
(\delta_1^{x_i})^g
=
\delta_1^{x_i g}
=
\delta_1^{g(g^{-1}x_i g)}
=
\delta_2^{g^{-1}x_i g}.
\]
Since \(N \lhd G\), we have \(g^{-1}x_i g \in N\), and therefore
\[
(\delta_1^{x_i})^g \in \Delta_2.
\]
Thus
\[
\Delta_1^g \subseteq \Delta_2.
\]
Applying the same argument to \(g^{-1}\), we get equality:
\[
\Delta_1^g=\Delta_2.
\]
Hence \(G\) permutes the \(N\)-orbits transitively, so all \(N\)-orbits have the
same size.
\end{proof}


\begin{corollary}
Every primitive permutation group is quasiprimitive.
\end{corollary}

\begin{proof}
Let \(G\) be primitive, and let \(1\neq N\lhd G\). By the previous lemma, the
\(N\)-orbits form a \(G\)-invariant block system. Since \(G\) is primitive, this block
system is either trivial or consists of singleton blocks.

If the \(N\)-orbits are singleton blocks, then every element of \(N\) fixes every
point of \(\Omega\), so \(N=1\), a contradiction. Hence \(N\) is transitive.
\end{proof}

\section{Regular Nonabelian Minimal Normal Subgroups}

\begin{definition}[Holomorph simple type]
    Let \(T\) be a nonabelian simple group. A quasiprimitive permutation group
    \(G \leq \operatorname{Sym}(\Omega)\) is said to be of \textbf{holomorph simple
    type}, or of type \(\mathrm{HS}\), if \(G\) has exactly two minimal normal
    subgroups
    \[
        M
        \quad\text{and}\quad
        N,
    \]
    such that
    \[
        M\cong N\cong T,
    \]
    and both \(M\) and \(N\) act regularly on \(\Omega\).
    
    Equivalently, after identifying \(\Omega\) with the underlying set of \(M\) via
    the regular action of \(M\), we have
    \[
        M:\operatorname{Inn}(M)
        \leq
        G
        \leq
        M:\operatorname{Aut}(M).
    \]
    Here \(M:\operatorname{Aut}(M)\) is the holomorph of \(M\), acting on \(M\)
    by right multiplication and automorphisms.
    \end{definition}

\begin{theorem}[Rank bound for holomorph simple type]
Let \(G\) be of holomorph simple type. Then
\[
\operatorname{rank}(G)\geq 4.
\]
\end{theorem}

\begin{proof}
Identify \(\Omega\) with \(T\), and take the base point to be \(1\in T\). Then
\[
G_1 \leq \operatorname{Aut}(T).
\]
The \(G_1\)-orbits on \(T\) are contained in automorphism orbits, and automorphisms
preserve element orders.

Since \(T\) is nonabelian simple, \(|T|\) has at least three distinct prime divisors.
Indeed, \(T\) is not solvable; groups of odd order are solvable, and groups of order
\(p^a q^b\) are solvable by Burnside's \(p^a q^b\)-theorem. Hence there exist
elements
\[
x,y,z \in T
\]
of three distinct prime orders.

The four elements
\[
1,\quad x,\quad y,\quad z
\]
lie in four distinct \(G_1\)-orbits, because their orders are distinct. Therefore
\[
\operatorname{rank}(G)\geq 4.
\]
\end{proof}

\begin{example}[Sharpness of the bound]
Let \(T=A_5\), and let
\[
G=T\rtimes \operatorname{Aut}(T) \cong A_5 \rtimes S_5.
\]
In the holomorph action of \(G\) on \(T\), the point stabilizer is
\[
G_1=\operatorname{Aut}(A_5)\cong S_5.
\]
The conjugation action of \(S_5\) on \(A_5\) has exactly four orbits:
\[
\{1\},
\]
the elements of order \(2\), the elements of order \(3\), and the elements of order
\(5\). Therefore
\[
\operatorname{rank}(G)=4.
\]
Thus the bound \(\operatorname{rank}(G)\geq 4\) is best possible.
\end{example}

\begin{proposition}[Compound regular case]
    Let \(N=T^k\), where \(k\geq 2\) and \(T\) is a finite nonabelian simple group.
    Suppose that \(N\) is a regular normal subgroup of a transitive permutation group
    \(G\). Then
    \[
        \operatorname{rank}(G)\geq 10.
    \]
    \end{proposition}

    \begin{proof}
        Fix \(\omega\in \Omega\). Since \(N\) is regular on \(\Omega\), we may identify
        \(\Omega\) with \(N\) by
        \[
            n \longleftrightarrow \omega^n.
        \]
        Under this identification, \(\omega\) corresponds to the identity element
        \[
            \mathbf{1}=(1,\ldots,1)\in N.
        \]
        
        Because \(N\lhd G\), the point stabilizer \(G_\omega\) normalizes \(N\). Hence
        conjugation gives a homomorphism
        \[
            G_\omega \longrightarrow \operatorname{Aut}(N),
            \qquad
            g\longmapsto (n\mapsto g^{-1}ng).
        \]
        This homomorphism is injective. Indeed, if \(g\in G_\omega\) centralizes \(N\),
        then for every \(n\in N\),
        \[
            (\omega^n)^g
            =
            \omega^{ng}
            =
            \omega^{g(g^{-1}ng)}
            =
            (\omega^g)^{g^{-1}ng}
            =
            \omega^n.
        \]
        Thus \(g\) fixes every point of \(\Omega\), so \(g=1\). Therefore we may regard
        \[
            G_\omega \leq \operatorname{Aut}(N).
        \]
        
        Write
        \[
            N=T_1\times T_2\times \cdots \times T_k,
        \]
        where each \(T_i\cong T\). The subgroups \(T_1,\ldots,T_k\) are precisely the
        minimal normal subgroups of \(N\). Hence every automorphism of \(N\) permutes
        the direct factors \(T_i\). Consequently, if
        \[
            x=(x_1,\ldots,x_k)\in N,
        \]
        then the number
        \[
            \operatorname{supp}(x)
            =
            |\{i\mid x_i\neq 1\}|
        \]
        is invariant under \(\operatorname{Aut}(N)\), and hence is invariant under
        \(G_\omega\). Also, since elements of \(G_\omega\) act as automorphisms of \(N\),
        they preserve element orders. Therefore the pair
        \[
            \bigl(\operatorname{supp}(x), |x|\bigr)
        \]
        is constant on each \(G_\omega\)-orbit on \(N\).
        
        Since \(T\) is finite nonabelian simple, \(|T|\) has at least three distinct prime
        divisors. Indeed, otherwise \(|T|\) would have the form \(p^a\) or \(p^aq^b\), and
        \(T\) would be solvable, contradicting the nonabelian simplicity of \(T\). Choose
        three distinct primes
        \[
            p,\ q,\ r \mid |T|,
        \]
        and choose elements
        \[
            a,b,c\in T
        \]
        such that
        \[
            |a|=p,\qquad |b|=q,\qquad |c|=r.
        \]
        
        Now consider the following elements of \(N=T_1\times\cdots\times T_k\):
        \[
            \mathbf{1}=(1,\ldots,1),
        \]
        \[
            (a,1,\ldots,1),\qquad
            (b,1,\ldots,1),\qquad
            (c,1,\ldots,1),
        \]
        \[
            (a,a,1,\ldots,1),\qquad
            (b,b,1,\ldots,1),\qquad
            (c,c,1,\ldots,1),
        \]
        and
        \[
            (a,b,1,\ldots,1),\qquad
            (a,c,1,\ldots,1),\qquad
            (b,c,1,\ldots,1).
        \]
        These elements have pairwise distinct values of the invariant
        \[
            \bigl(\operatorname{supp}(x), |x|\bigr).
        \]
        Indeed, the corresponding pairs are
        \[
            (0,1),
        \]
        \[
            (1,p),\qquad (1,q),\qquad (1,r),
        \]
        \[
            (2,p),\qquad (2,q),\qquad (2,r),
        \]
        and
        \[
            (2,pq),\qquad (2,pr),\qquad (2,qr).
        \]
        These ten pairs are pairwise distinct. Hence the ten elements listed above lie in ten
        distinct \(G_\omega\)-orbits on \(N\).
        
        Therefore \(G_\omega\) has at least ten orbits on \(\Omega\). Since the rank of a
        transitive permutation group is the number of point-stabilizer orbits, we obtain
        \[
            \operatorname{rank}(G)\geq 10.
        \]
        \end{proof} 



\begin{corollary}
A quasiprimitive group of rank at most \(4\) cannot have a regular minimal normal
subgroup \(T^k\) with \(k\geq 2\) and \(T\) nonabelian simple.
\end{corollary}

\section{A Useful Obstruction in the Two-Transitive Case}

\begin{lemma}
Let \(G\leq \operatorname{Sym}(\Omega)\) be \(2\)-transitive. Let
\[
N=T^k \lhd G,
\]
where \(k\geq 2\) and \(T\) is nonabelian simple. Then \(N\) has no proper normal
subgroup that is regular on \(\Omega\).
\end{lemma}

\begin{proof}
Suppose, for a contradiction, that \(M\lhd N\) is regular on \(\Omega\), with
\(M<N\). Since \(M\) is regular,
\[
|\Omega|=|M|.
\]
Because \(N=T^k\) is a direct product of simple groups, every normal subgroup of
\(N\) is a product of some of the direct factors. Hence
\[
|M|=|T|^\ell
\]
for some integer \(0<\ell<k\).

Fix \(\omega\in \Omega\). Since \(M\) is regular, \(M_\omega=1\). Also
\[
N=MN_\omega,
\]
so
\[
|N_\omega|=\frac{|N|}{|M|}
=
|T|^{k-\ell}.
\]

Because \(G\) is \(2\)-transitive, \(N_\omega\) is half-transitive on
\(\Omega\setminus\{\omega\}\). Let \(m\) be the common size of its orbits on
\(\Omega\setminus\{\omega\}\). Then
\[
m \mid |\Omega|-1
\]
and
\[
m\mid |N_\omega|.
\]
Hence
\[
m \mid \gcd(|\Omega|-1,|N_\omega|)
=
\gcd(|T|^\ell-1,|T|^{k-\ell})
=
1.
\]
Therefore \(m=1\).

Thus \(N_\omega\) fixes every point of \(\Omega\setminus\{\omega\}\), and it also
fixes \(\omega\). Hence \(N_\omega=1\), contradicting \(N_\omega\neq 1\), since
\(M<N\) and \(N=MN_\omega\). Therefore no such proper regular normal subgroup
\(M\) exists.
\end{proof}

\section{Product Action Type}

\begin{definition}[Product action]
Let \(T\leq \operatorname{Sym}(\Delta)\) be transitive, and let \(k\geq 2\). The
wreath product
\[
T \wr S_k = T^k \rtimes S_k
\]
acts on
\[
\Omega=\Delta^k
\]
by product action. Explicitly, for
\[
g=(t_1,\ldots,t_k)\pi \in T^k\rtimes S_k
\]
and
\[
\delta=(\delta_1,\ldots,\delta_k)\in \Delta^k,
\]
we define
\[
\delta^g
=
(\delta_{1\pi^{-1}}^{t_{1\pi^{-1}}},
 \ldots,
 \delta_{k\pi^{-1}}^{t_{k\pi^{-1}}}).
\]
\end{definition}

\begin{proposition}[Product action reduction]
Let \(G\leq \operatorname{Sym}(\Omega)\) be primitive, and suppose
\[
N=T_1\times \cdots \times T_k \lhd G
\]
is a transitive nonregular minimal normal subgroup, where
\[
T_1\cong \cdots \cong T_k\cong T
\]
and \(T\) is nonabelian simple.

Fix \(\omega\in \Omega\), and assume that the projection of \(N_\omega\) to each
factor \(T_i\) is a proper subgroup \(H_i<T_i\). Then
\[
N_\omega=H_1\times \cdots \times H_k,
\]
and
\[
\Omega
\cong
[T_1:H_1]\times \cdots \times [T_k:H_k].
\]
Under this identification, \(N\) acts coordinatewise by right multiplication, and
\(G\) embeds into a wreath product in product action.
\end{proposition}

\begin{proof}
Since \(N\) is transitive,
\[
\Omega \cong [N:N_\omega].
\]
Write
\[
N=T_1\times \cdots \times T_k.
\]
Let \(H_i\) be the image of \(N_\omega\) under projection onto \(T_i\). Then
\[
N_\omega \leq H_1\times \cdots \times H_k.
\]
Set
\[
K=H_1\times \cdots \times H_k.
\]
Since the factors \(T_i\) are permuted transitively by \(G_\omega\), the subgroup
\(K\) is normalized by \(G_\omega\). Moreover, by hypothesis \(H_i<T_i\), so
\[
K<N.
\]

Because \(G\) is primitive, \(G_\omega\) is a maximal subgroup of \(G\). Hence
\[
G_\omega \leq KG_\omega \leq G
\]
implies either
\[
KG_\omega=G_\omega
\]
or
\[
KG_\omega=G.
\]
The second possibility is impossible. Indeed, if \(KG_\omega=G\), then
\[
N=(KG_\omega)\cap N
=
K(G_\omega\cap N)
=
KN_\omega
=
K,
\]
contradicting \(K<N\). Hence
\[
KG_\omega=G_\omega,
\]
so \(K\leq G_\omega\). Since also \(K\leq N\), we get
\[
K\leq N\cap G_\omega=N_\omega.
\]
Thus
\[
N_\omega=K=H_1\times \cdots \times H_k.
\]

Therefore
\[
\Omega
\cong
[N:N_\omega]
=
[T_1\times \cdots \times T_k : H_1\times \cdots \times H_k]
\]
and hence
\[
\Omega
\cong
[T_1:H_1]\times \cdots \times [T_k:H_k].
\]
This is precisely the product action model.
\end{proof}

\section{Ranks in Product Action}

\begin{definition}[Weight]
Let
\[
\Delta=\{0,1,\ldots,n-1\}
\]
and let
\[
\Omega=\Delta^k.
\]
For
\[
u=(\delta_1,\ldots,\delta_k)\in \Omega,
\]
define the \textbf{weight} of \(u\) by
\[
\operatorname{wt}(u)
=
|\{i\mid \delta_i\neq 0\}|.
\]
Thus \(\operatorname{wt}(u)\) is the number of nonzero coordinates of \(u\).
\end{definition}

\begin{lemma}[Rank in full product action]
Let \(T\leq \operatorname{Sym}(\Delta)\) be transitive, and let
\[
G=T^k\rtimes S_k
\]
act on
\[
\Omega=\Delta^k
\]
in product action. Fix \(0\in \Delta\), and set
\[
\omega=(0,\ldots,0).
\]
Then
\[
\operatorname{rank}(G)\geq k+1.
\]
Moreover,
\[
\operatorname{rank}(G)=k+1
\]
if and only if \(T\) is \(2\)-transitive on \(\Delta\).
\end{lemma}

\begin{proof}
Let \(H=T_0\). Then
\[
G_\omega=H^k\rtimes S_k.
\]
The group \(G_\omega\) preserves the weight of a point in \(\Omega\). Therefore the
sets
\[
\Omega_j=\{u\in \Omega\mid \operatorname{wt}(u)=j\},
\qquad 0\leq j\leq k,
\]
are unions of \(G_\omega\)-orbits. Hence
\[
\operatorname{rank}(G)\geq k+1.
\]

Suppose first that \(T\) is \(2\)-transitive. Then \(H=T_0\) is transitive on
\[
\Delta\setminus\{0\}.
\]
Thus, if two points of \(\Omega\) have the same weight, then a permutation in
\(S_k\) moves the nonzero coordinates of one point to the nonzero coordinates of
the other, and elements of \(H\) adjust the values in those coordinates. Hence points
of the same weight lie in the same \(G_\omega\)-orbit. Therefore the \(G_\omega\)-orbits
are exactly the \(k+1\) weight classes, and
\[
\operatorname{rank}(G)=k+1.
\]

Conversely, suppose \(T\) is not \(2\)-transitive. Then \(H=T_0\) has at least two
orbits on \(\Delta\setminus\{0\}\). Choose representatives
\[
a,b\in \Delta\setminus\{0\}
\]
from two distinct \(H\)-orbits. Then
\[
(a,0,\ldots,0)
\quad\text{and}\quad
(b,0,\ldots,0)
\]
have the same weight, but they are not \(G_\omega\)-conjugate. Hence the weight-one
class splits into at least two \(G_\omega\)-orbits, so
\[
\operatorname{rank}(G)>k+1.
\]
Thus equality holds if and only if \(T\) is \(2\)-transitive.
\end{proof}

\begin{example}[The case \(k=2\)]
Let
\[
G=(T\times T)\rtimes S_2
\]
act on
\[
\Omega=\Delta^2.
\]
Then
\[
\operatorname{rank}(G)\geq 3,
\]
and
\[
\operatorname{rank}(G)=3
\]
if and only if \(T\) is \(2\)-transitive on \(\Delta\).

For example, take \(T=A_5\) acting naturally on
\[
\Delta=\{0,1,2,3,4\}.
\]
Then \(T\) is \(2\)-transitive, and
\[
G=(A_5\times A_5)\rtimes S_2
\]
has rank \(3\) on \(\Delta^2\).

Fix
\[
\omega=(0,0).
\]
Then
\[
G_\omega=(A_4\times A_4)\rtimes S_2.
\]
The nontrivial \(G_\omega\)-orbits on \(\Omega\setminus\{\omega\}\) are
\[
\{(0,i),(i,0)\mid 1\leq i\leq 4\}
\]
and
\[
\{(i,j)\mid 1\leq i,j\leq 4\}.
\]
Together with \(\{\omega\}\), these give exactly three suborbits.
\end{example}

\begin{exercise}
Let
\[
G=(T\times T\times T)\rtimes S_3
\]
act on
\[
\Omega=\Delta^3
\]
in product action. Prove that
\[
\operatorname{rank}(G)\geq 4,
\]
and that
\[
\operatorname{rank}(G)=4
\]
if and only if \(T\) is \(2\)-transitive on \(\Delta\).
\end{exercise}

\begin{proof}
This is the case \(k=3\) of the previous lemma.

Explicitly, fix
\[
\omega=(0,0,0).
\]
The point stabilizer preserves the weight
\[
\operatorname{wt}(\delta_1,\delta_2,\delta_3)
=
|\{i\mid \delta_i\neq 0\}|.
\]
Therefore there are at least four suborbits, corresponding to weights
\[
0,1,2,3.
\]
Hence
\[
\operatorname{rank}(G)\geq 4.
\]

If \(T\) is \(2\)-transitive, then \(T_0\) is transitive on
\(\Delta\setminus\{0\}\), so points of the same weight are \(G_\omega\)-conjugate.
Thus the four weight classes are exactly the four suborbits, and
\[
\operatorname{rank}(G)=4.
\]

Conversely, if \(T\) is not \(2\)-transitive, then \(T_0\) has at least two orbits on
\(\Delta\setminus\{0\}\). Hence the weight-one class splits into at least two
\(G_\omega\)-orbits. Therefore
\[
\operatorname{rank}(G)>4.
\]
So equality holds if and only if \(T\) is \(2\)-transitive.
\end{proof}

\section{Summary for Small Rank}

\begin{topic}
For the small-rank problem, the following constraints are especially useful.

\begin{enumerate}[label=\textup{(\arabic*)}]
    \item \(\operatorname{rank}(G)=2\) if and only if \(G\) is \(2\)-transitive.
    By Burnside's theorem, a finite \(2\)-transitive group is either affine type
    or almost simple type.

    \item If \(G\) is primitive, then \(G\) is quasiprimitive.

    \item If \(G\) is of holomorph simple type, then
    \[
    \operatorname{rank}(G)\geq 4.
    \]
    This bound is sharp, as shown by the holomorph action of
    \(A_5\rtimes S_5\) on \(A_5\).

    \item If \(G\) has a regular normal subgroup \(T^k\), where
    \(k\geq 2\) and \(T\) is nonabelian simple, then
    \[
    \operatorname{rank}(G)\geq 10.
    \]

    \item In product action
    \[
    G=T^k\rtimes S_k
    \quad\text{on}\quad
    \Omega=\Delta^k,
    \]
    one always has
    \[
    \operatorname{rank}(G)\geq k+1.
    \]
    Equality holds precisely when \(T\) is \(2\)-transitive on \(\Delta\).
\end{enumerate}
\end{topic}